\section{Prototipazione Low-Fi e Mago di Oz}

\subsection{Paper Sketch Finali}
% Scrivere qui la descrizione dei paper sketch finali e di come sono stati digitalizzati/organizzati.

\begin{figure}[H]
    \centering
    % \includegraphics[width=\textwidth]{img/paper_sketches.png}
    \framebox[\textwidth]{\vbox to 8cm{\vfill\centering [Inserire qui i Paper Sketch finali caricate su Figma]\vfill}}
    \caption{Paper sketch finali del sistema.}
    \label{fig:paper_sketches}
\end{figure}

\subsection{Tecnica del Mago di Oz Adottata}
% Descrivere dettagliatamente la metodologia del Mago di Oz utilizzata per valutare l'usabilità del prototipo low-fi rispetto ai task ed agli scenari considerati.

\subsection{Risultati di Usabilità ed Iterazioni}
% Riassumere i risultati ottenuti dai test ed evidenziare tutte le eventuali iterazioni/modifiche che si sono rese necessarie prima di arrivare al prototipo finale.
